\section*{Chapter 5: Conclusion and Further Work}
\phantomsection \setcounter{section}{5} \setcounter{subsection}{0}
\addcontentsline{toc}{section}{\textit{Chapter 5: Conclusion and Further Work}}

\subsection{Conclusion}

This project set out to build a hybrid deep learning model for automated epilepsy detection from EEG signals and to evaluate it under a stricter protocol than the initial implementation. The final results support the usefulness of that design on the Bonn benchmark.

The proposed 1D-CNN + Bi-LSTM + Self-Attention architecture achieves 99.67\% on binary classification, 97.69\% on three-class, and 84.51\% on the five-class task --- all evaluated under strict recording-level 10-Fold cross-validation with no data leakage. On Bonn five-class classification, the reported mean accuracy is higher than the 67.24\% reported by Huang et al.\ \cite{huang2025}, although the preprocessing pipelines differ and the comparison should therefore be treated as indicative rather than exact.

The ablation study suggests that the three components contribute complementary strengths. Paired t-tests show the hybrid model significantly outperforms both the pure CNN ($p = 0.001$) and pure Bi-LSTM ($p = 0.001$) on five-class classification. The binary task, by contrast, shows no significant differences in the reported pairwise tests, which is consistent with a ceiling effect in this particular Bonn setup.

Attention weight visualisation highlights high-amplitude seizure regions and offers a qualitative view of model behaviour. It is useful for discussion, but it should not be treated as a clinically validated explanation.

\subsection{Reflection}

Looking back, the most useful lesson from this project concerned evaluation methodology more than model design.

Early in the project, my initial preprocessing pipeline produced near-perfect accuracy --- over 99.9\% on binary classification. The numbers looked appealing, but my supervisor rightly questioned them. After investigation, I found that overlapping sliding-window segments from the same recording were leaking across the train/validation split. Fixing this dropped accuracy by about 0.3\% on binary and 2\% on five-class. The revised numbers are less dramatic, but they are more credible. This part of the project taught me that evaluation design matters at least as much as model design.

On the technical side, migrating the codebase from TensorFlow to PyTorch mid-project was not planned but turned out to be necessary --- TensorFlow 2.11 dropped native Windows GPU support, and under my local setup training on CPU was far slower than PyTorch with CUDA. The migration forced me to understand both frameworks at a deeper level than I would have otherwise.

I also want to be transparent about the role of AI tools. I used Cursor as a support tool for code scaffolding, debugging, and some report drafting. All generated code and text were reviewed and revised by me. The main project decisions --- including the model choice, the leakage correction, and the experimental design --- remained my responsibility.

\subsection{Further Work}

Several directions could extend this work:

\begin{itemize}
    \item \textbf{Larger, multi-channel datasets.} The Bonn dataset has only 500 single-channel recordings and is relatively clean. Validating on noisier, multi-channel datasets like CHB-MIT \cite{shoeb2010} or the Temple University Hospital EEG Corpus \cite{obeid2016} would test whether the model generalises to real clinical conditions.

    \item \textbf{Hyperparameter sensitivity analysis.} The current hyperparameters (window size 1024, dropout 0.4, LSTM units 128/64) were chosen based on literature defaults and preliminary experiments. A systematic grid search or Bayesian optimisation over these parameters could improve performance or confirm robustness.

    \item \textbf{Cross-dataset transfer learning.} Pre-training on a large, diverse EEG corpus and fine-tuning on a smaller target dataset could improve generalisation in data-scarce clinical settings.

    \item \textbf{Real-time inference.} The current system processes pre-recorded segments. Deploying it for real-time EEG monitoring would require model compression (pruning, quantisation) and optimisation for edge hardware.

    \item \textbf{Patient-specific adaptation.} EEG patterns vary substantially between individuals. Developing mechanisms for the model to adapt to a specific patient's baseline could improve detection accuracy in personalised settings.
\end{itemize}
